\section{A scalar growth limit with the exact critical root}
\label{gc:section}

The inert comparison gives an actual generation limit, and the preceding continuous-time comparison identifies its critical parameter.

\begin{theorem}[Exact inert-parent characteristic equation]
\label{gc:root}
For $s>0$, let
\[
 r_\infty(s)=\lim_{m\to\infty} A_m^\infty(s)^{1/m}.
\]
Then $r_\infty$ is positive, continuous, and strictly decreasing, and
\begin{equation}\label{gc:crossing}
 \boxed{r_\infty(\mu)=1,\qquad
 r_\infty(s)>1\ (0<s<\mu),\qquad
 r_\infty(s)<1\ (s>\mu).}
\end{equation}
In particular the original exponent is the unique solution of this
scalar equation. For every $m\ge1$,
\begin{equation}\label{gc:critical-bound}
 A_m^\infty(\mu)\le 12(\mu+3)<43.
\end{equation}
\end{theorem}
\begin{proof}
Tonelli's theorem gives the extended-value identity
\[
 1+\sum_{m\ge1}A_m^\infty(s)
 =s\int_0^\infty e^{-st}F_\infty(t)\,dt.
\]
By Theorem~\ref{gs:inert-mu}, this is finite for $s>\mu$ and infinite
for $s<\mu$. Hence $r_\infty(s)\le1$ above $\mu$ and
$r_\infty(s)\ge1$ below $\mu$.

Each $A_m^\infty$ is a Laplace transform of a positive measure, so
$m^{-1}\log A_m^\infty(s)$ is convex. Its pointwise limit
$g(s)=\log r_\infty(s)$ exists by Theorem~\ref{growth:inert-limit}
and has the finite bounds
\[
 -\log\bigl(2(s+2)\bigr)\le g(s)\le-\log s.
\]
Thus $g$ is finite, convex, continuous, and nonincreasing on
$(0,\infty)$. If it took the same value at two distinct points,
convexity and monotonicity would force it to be constant thereafter.
This contradicts $g(s)\le-\log s\to-\infty$. It is strictly decreasing.
The weak inequalities on either side of $\mu$, together with continuity
and strict decrease, prove~\eqref{gc:crossing}.
Equation~\eqref{growth:uniform-inert-upper} at $s=\mu$ gives
\eqref{gc:critical-bound}; the numerical inequality uses the imported
upper bound $\mu\le1133/2000$.
\end{proof}

\begin{corollary}[Complete scalar lower certificates]
\label{gc:scalar-certificates}
For every $s>0$,
\begin{equation}\label{gc:certificate-equivalence}
 \boxed{s<\mu\quad\Longleftrightarrow\quad
 A_m^\infty(s)>12(s+3)\text{ for some finite }m.}
\end{equation}
At rational $s$, finite positive rational-history approximations provide
a complete semidecision procedure for $s<\mu$. The same statement
holds for a positive algebraic $s$ specified exactly: each finite
comparison is then an exact algebraic-number sign calculation.
\end{corollary}
\begin{proof}
Combine~\eqref{growth:sup-formula} with~\eqref{gc:crossing}.
When the inequality holds, its strict margin is eventually retained
by finite positive history approximations. Conversely every successful
finite lower approximation is bounded by the exact quantity and is
therefore sound. All its weights are rational functions of $s$ with
positive denominators for $s>0$.
\end{proof}

For all sufficiently large $m$, the equation
\[
 \frac{A_m^\infty(s_m)}{12(s_m+3)}=1
\]
has a unique positive solution and $s_m\le\mu$, with $s_m\to\mu$.
Indeed, for each fixed $s<\mu$, the left side at $s$ eventually exceeds
one by the actual generation limit; at $\mu$ it is at most one and at
infinity it tends to zero. Each individual left side is continuous and
strictly decreasing. The running maxima of these roots therefore
increase to $\mu$. No effective outer error or monotonicity of the
individual $s_m$ is asserted.

